\begin{lemma}[Structural equivalence is an equivalence relation]
The relation $\mathsf{PowerPresentationEq}$ is reflexive, symmetric, and
transitive on $\mathsf{PowerQ}(Q)$.
\end{lemma}
\begin{proof}
The induction and all case distinctions below use the two constructors of the
tagged inductive carrier \noderef{PowerQ}.  We unfold the recursive relation
from \noderef{PowerPresentationEq}; its atom--node and node--atom clauses are
both the proposition \(\mathsf{False}\).

For reflexivity, induct on \(x\in\mathsf{PowerQ}(Q)\).  If
\(x=\mathsf{atom}(q)\), the atom--atom clause asks for \(q=q\), which is
given by reflexivity of equality.  If
\(x=\mathsf{node}(\iota,h,f)\), the relation has two conjuncts.  For each
\(i\in\iota\), choose the same index \(i\) on the other side; the induction
hypothesis gives
\(\mathsf{PowerPresentationEq}(f(i),f(i))\).  This proves the forward
conjunct \(\forall i\,\exists j\), and choosing the same index in the
reverse direction proves the backward conjunct \(\forall j\,\exists i\).
Thus every presentation is related to itself.  The nonemptiness witness
\(h\) is part of the node constructor supplied by \noderef{PowerQ}, although
it is not needed for these two matching clauses.

For symmetry, perform the same simultaneous structural induction on the two
presentations occurring in a relation.  In the atom--atom case,
\(q=q'\) reverses to \(q'=q\).  If exactly one presentation is an atom, the
assumed relation is \(\mathsf{False}\), so there is no case to prove; these
are precisely the atom--node and node--atom clauses of
\noderef{PowerPresentationEq}.  In the node--node case, write the assumed
relation as
\[
 H_{\to}:\forall i\in\iota,\ \exists j\in\kappa,\ \mathsf{PowerPresentationEq}(f(i),g(j)),
 \qquad
 H_{\leftarrow}:\forall j\in\kappa,\ \exists i\in\iota,\ \mathsf{PowerPresentationEq}(f(i),g(j)).
\]
Given \(j\in\kappa\), choose the witness \(i\) supplied by
\(H_{\leftarrow}\); the induction hypothesis reverses its matching relation
to \(\mathsf{PowerPresentationEq}(g(j),f(i))\).  This proves the forward
clause for the reversed pair.  Given \(i\in\iota\), use \(H_{\to}\) and
the induction hypothesis in the same way to prove the reversed backward
clause.  Hence the relation is symmetric.

For transitivity, let \(x,y,z\in\mathsf{PowerQ}(Q)\) and assume
\(\mathsf{PowerPresentationEq}(x,y)\) and
\(\mathsf{PowerPresentationEq}(y,z)\).  Use simultaneous structural
induction on these three presentations.  If all three are atoms, write the
two hypotheses as \(q=q'\) and \(q'=q''\); their composite gives
\(q=q''\).  Every mixed-tag arrangement has a false premise: whenever the
first and middle tags differ, the first relation is one of the two mixed
\(\mathsf{False}\) clauses, and whenever the middle and last tags differ, the
second relation is such a clause.  Thus the arrangements with one or two
atom--node changes, including atom--node and node--atom in either order,
are impossible and require no further construction.

It remains to consider three nodes
\(x=\mathsf{node}(\iota,h,f)\),
\(y=\mathsf{node}(\kappa,h',g)\), and
\(z=\mathsf{node}(\lambda,h'',k)\).  Write the first relation as its
forward and backward matching clauses
\[
 H_{\to}:\forall i\in\iota,\ \exists j\in\kappa,\ \mathsf{PowerPresentationEq}(f(i),g(j)),
 \qquad
 H_{\leftarrow}:\forall j\in\kappa,\ \exists i\in\iota,\ \mathsf{PowerPresentationEq}(f(i),g(j)),
\]
and write the second relation similarly as
\[
 K_{\to}:\forall j\in\kappa,\ \exists \ell\in\lambda,\ \mathsf{PowerPresentationEq}(g(j),k(\ell)),
 \qquad
 K_{\leftarrow}:\forall \ell\in\lambda,\ \exists j\in\kappa,\ \mathsf{PowerPresentationEq}(g(j),k(\ell)).
\]
For the forward clause of the composite, fix \(i\in\iota\).  Choose
\(j\in\kappa\) and a witness
\(e:\mathsf{PowerPresentationEq}(f(i),g(j))\) from \(H_{\to}\).  Choose
\(\ell\in\lambda\) and a witness
\(d:\mathsf{PowerPresentationEq}(g(j),k(\ell))\) from \(K_{\to}\).
The induction hypothesis for these three children composes \(e\) and \(d\),
giving \(\mathsf{PowerPresentationEq}(f(i),k(\ell))\).  This is the required
\(\ell\)-witness.  For the backward clause, fix \(\ell\in\lambda\).  Use
\(K_{\leftarrow}\) to choose \(j\in\kappa\) and its matching witness, then
use \(H_{\leftarrow}\) to choose \(i\in\iota\) matching that middle child.
The induction hypothesis again composes the two witnesses and gives
\(\mathsf{PowerPresentationEq}(f(i),k(\ell))\).  These two constructions
establish the node--node relation from \(x\) to \(z\), so the relation is
transitive and therefore an equivalence relation.
\end{proof}
